% Generated draft: requires mathematical and source review.
\documentclass[11pt]{article}
\usepackage[margin=25mm]{geometry}
\usepackage{fontspec}
\setmainfont{Latin Modern Roman}
\usepackage{amsmath,amssymb,mathtools,unicode-math}
\setmathfont{Latin Modern Math}
\usepackage{enumitem,needspace,xurl,hyperref}
\setlength{\parindent}{0pt}
\setlength{\parskip}{5pt}
\setlength{\emergencystretch}{3em}
\newcommand{\N}{\mathbb{N}}
\newcommand{\Z}{\mathbb{Z}}
\newcommand{\Q}{\mathbb{Q}}
\newcommand{\R}{\mathbb{R}}
\newcommand{\C}{\mathbb{C}}
\newcommand{\F}{\mathbb{F}}
\newcommand{\A}{\mathbb{A}}
\newcommand{\PP}{\mathbb{P}}
\newcommand{\E}{\mathbb{E}}
\newcommand{\eps}{\varepsilon}
\newcommand{\dd}{\,\mathrm{d}}
\begin{document}
\section*{UnsolvedMath ID: 10400142}
\textbf{Source problem code:} catalogue entry
\subsection{Definitions and mathematical statement}
For orientation of terminology, a scalar modular form of integral weight \(k\) for a subgroup \(\Gamma\subset\mathrm{SL}_2(\Z)\), in the simplest trivial-multiplier convention, is a holomorphic function on \(\mathbb H=\{\tau\in\C:\operatorname{Im}\tau>0\}\) satisfying
\[
 f\!\left(\frac{a\tau+b}{c\tau+d}\right)
      =(c\tau+d)^k f(\tau)
 \quad\left(\begin{pmatrix}a&b\\c&d\end{pmatrix}\in\Gamma\right),
\]
with the prescribed cusp holomorphy conditions. Other weights, multipliers and vector-valued forms require their own conventions; the source selects none here.

J. Roberts's Problem 7.27 asks to explain the appearance of modular forms in the Witten invariants. Write \(\tau_r^G(M)\) for a Witten-type quantum invariant only as descriptive notation: the record does not specify the group \(G\), level convention \(r\), normalization, manifold class, or analytic series to which a modular transformation law is attached. It gives no particular modular form, weight, group, or identity to prove.

The requested explanation therefore remains a comparison program: identify the relevant objects and the mechanism relating them in the intended examples. It must not be replaced by the assertion that every quantum invariant is itself a modular form. The supplied partial-progress summary points to modularity phenomena for quantum invariants at roots of unity and work of Zagier, Lawrence and others, but supplies neither exact examples nor complete primary references.
\subsection{Short English statement}
Explain why modular forms occur in the study of Witten invariants. The source does not select the invariants, manifolds, modular objects, or transformation formulas. Its supplied assessment reports related work and partial progress rather than a precise resolved theorem.
\subsection{Sources}
\begin{itemize}
\item Local UnsolvedMath registry, permanent ID \(10400142\), exact source code \texttt{catalogue entry}, Problem 7.27, attributed to J. Roberts. The supplied title is truncated after ``(J.''
\item Ohtsuki, \emph{Problems on invariants of knots and 3-manifolds} (2002), Problem 7.27, PDF page 118: \url{https://msp.org/gtm/2002/04/gtm-2002-04-024s.pdf}. Bold printed headings distinguish 231 numbered items from surrounding references.
\item Extraction pdf-font-structure, accessed 2026-07-29. Source evidence NEEDS\_REVIEW presents an open problem and requires a current primary-literature audit. Registry status partially\_solved and classification PARTIAL-PROGRESS mention roots-of-unity modularity and work of Zagier and Lawrence, without complete citations or a precise resolved scope.
\item Dataset attribution: Ulam AI / UnsolvedMath Contributors, CC BY 4.0. Source rights NEEDS\_REVIEW: publicly accessible author/publisher copy, redistribution terms requiring review.
\end{itemize}
\end{document}
